\documentclass[11pt]{article}
\usepackage{fontspec}
\usepackage[french,english]{babel}
\usepackage{amsmath}
\usepackage{amssymb}
\usepackage{amsthm}
\usepackage{mathrsfs}
\usepackage{tikz-cd}
\usepackage{enumitem}

\newcommand{\printedpage}[1]{\marginpar{\footnotesize p.~#1}}
\newcommand{\msp}[1]{\marginpar{\footnotesize [#1]}}
\newcommand{\Ch}{\underline{\mathrm{Ch}}}
\newcommand{\uHom}{\underline{\mathrm{Hom}}}
\newcommand{\uFl}{\underline{\mathrm{Fl}}}
\newcommand{\bDelta}{\boldsymbol{\Delta}}
\newcommand{\opp}{\mathrm{opp}}

\begin{document}

\printedpage{9}
% head: A. Grothendieck : Dérivateurs, Ch. VII, Édition des Universités de Montpellier II et Paris VII
% foot: Version 22 décembre 2001

\msp{page 13}

\section*{3\quad Paths. Category of paths of a given type $\tau$.}

Let $I$ be an oriented interval, and $X$ a category. We call a \textit{path of type $I$ in $X$} any functor $I \longrightarrow X$, and the \textit{category of paths of type $I$ in $X$} the category
\[
\Ch_I(X) = \uHom(I,X).
\]
Since $I$ is \textit{canonically} isomorphic to a standard interval category $I_0$, we can identify the paths of type $I$, with the of type $I_0$. The type $\tau$ of a path of length $n$ $(= \operatorname{card} I)$ in $X$ can therefore be described by a sequence of $n$ signs $+$ or $-$ % [sic] au de type -> à ceux de type, word(s) missing; [sic] le longueur -> de longueur, wrong preposition
\[
\tau_1,\ \tau_2,\ \dots\ ,\ \tau_n \qquad (\tau_i \in \{-1,+1\} \text{ for } 1 \leq i \leq n).
\]
Because of the fact that $I$ is generated by its transition arrows, a functor $c : I \longrightarrow X$ is determined when we know % [sic] pas ses flèches -> par ses flèches, misprint

\begin{enumerate}
\item[a)] The $c(i)$ $(i \in I)$, i.e. a sequence of $n+1$ objects of $X$:
\[
a_0,\ a_1,\ \dots\ ,\ a_n\ (\text{\footnotemark[5]});
\]
\footnotetext[5]{$a_0$ is called the \textit{source} of the path, $a_n$ its target.}

\item[b)] The $c(\underline{u}_i)$ (where the $\underline{u}_i$ are the transition arrows in $I$), i.e. for every $i \in [1,n]$ an arrow $u_i$ in $X$ between $a_{i-1}$ and $a_i$, going from $a_{i-1}$ to $a_i$ if $\underline{u}_i$ is direct $(\tau_i = +1)$, and from $a_i$ to $a_{i-1}$ if $\underline{u}_i$ is retrograde $(\tau_i = -1)$,
\end{enumerate}

and to any such datum a), b) there corresponds a functor. Thus we shall be able to identify the paths of type $I$ \msp{page 14} in $X$ with the diagrams of the specified type in $X$. And the morphisms between paths likewise correspond to the morphisms between diagrams $a$ and $a'$, an arrow $a \xrightarrow{\alpha} a'$ thus being a sequence of $n+1$ arrows
\[
\alpha_i : a_i \longrightarrow a'_i, \qquad i \in [0,n],
\]
satisfying the commutativity condition for the $n$ squares
\[
\begin{tikzcd}
a_{i-1} \arrow[r, "u_i"] \arrow[d, "\alpha_{i-1}"'] & a_i \arrow[d, "\alpha_i"] \\
a'_{i-1} \arrow[r, "u'_i"] & a'_i
\end{tikzcd}
\qquad\qquad \text{(case $\underline{u}_i$ direct)}
\]

\printedpage{10}
% head: A. Grothendieck : Dérivateurs, Ch. VII, Édition des Universités de Montpellier II et Paris VII
% foot: Version 22 décembre 2001
or
\[
\begin{tikzcd}
a_{i-1} \arrow[d, "\alpha_{i-1}"'] & a_i \arrow[l, "u_i"'] \arrow[d, "\alpha_i"] \\
a'_{i-1} & a'_i. \arrow[l, "u'_i"']
\end{tikzcd}
\qquad\qquad \text{(case $\underline{u}_i$ retrograde)}
\]
If we have an arrow between any two vertices $i$, $j$ in $I$, with $i < j$ and $j$ not the immediate successor of $i$, then $i$, $j$ belong to the same critical interval, and in $I$ we shall have % [sic] quelconque -> quelconques, agreement
\[
i \xrightarrow{\underline{u}_{i+1}} i+1 \xrightarrow{\underline{u}_{i+2}}\ \cdots\ j-1 \xrightarrow{\underline{u}_j} j \qquad \text{a chain of length } j-i \geq 2 % doubt resolved: PDF p. 12 (edition p. 10) clearly underlines u; the existing reading is retained.
\]
or
\[
i \longleftarrow i+1 \longleftarrow\ \ \cdots\ j-1 \longleftarrow j,
\]
and the corresponding arrow between $a_i$ and $a_j$ is obtained by composing the corresponding arrows $u_i$.

\msp{page 15}

\textbf{Examples.}

\begin{enumerate}
\item[1.] $I = \bDelta_0 = e$, $\Ch_I(X) \simeq X$. The paths of length 0 are identified with the objects $a$ of $X$. The unique element $\mathrm{id}_e$ of $\mathrm{Fl}(I)$ gives $\mathrm{id}_a$ in $X$.

\item[2.] $I = \bDelta_1^+$, $\Ch_I(X) \simeq \uFl(X)$. The paths of type $\bDelta_1^+$ are those of length 1 and \textit{direct,} i.e. given by an arrow
\[
a_0 \longrightarrow a_1
\]
whose source is also the source $a_0$ of the path. On the other hand, the paths of type $\bDelta_1^-$ are given by diagrams
\[
a_0 \longleftarrow a_1,
\]
i.e. by arrows whose target is regarded as the source of the path. Of course, the categories $\Ch_{\bDelta_1^+}(X)$ and $\Ch_{\bDelta_1^-}(X)$ are canonically isomorphic, but by an isomorphism which \textit{exchanges source and target of the path.} We shall therefore take good care not to confuse the two preceding paths.

\msp{page 16}

\item[3.] $\bDelta_2^+$, $\bDelta_2^-$, corresponding to the diagrams
\[
a_0 \longrightarrow a_1 \longrightarrow a_2
\]
resp.
\[
a_0 \longleftarrow a_1 \longleftarrow a_2,
\]
we can make a similar remark as for $\bDelta_1^+$, $\bDelta_1^-$.

\printedpage{11}
% head: A. Grothendieck : Dérivateurs, Ch. VII, Édition des Universités de Montpellier II et Paris VII
% foot: Version 22 décembre 2001

\item[4.] $\Phi_1$, $\Psi_1$, these are respectively the intervals
\[
\Phi_1 :\ 0 \xleftarrow{\underline{u}_0} 1 \xrightarrow{\underline{u}_1} 2 \qquad\qquad \text{or} \qquad\qquad
\begin{tikzcd}[row sep=small, column sep=small]
 & 0 \\
1 \arrow[ur] \arrow[dr] & \\
 & 2
\end{tikzcd}
\qquad , % doubt resolved: PDF p. 13 (edition p. 11) clearly underlines u in both labels; the existing reading is retained.
\]
and
\[
\Psi_1 :\ 0 \xrightarrow{\underline{u}_0} 1 \xleftarrow{\underline{u}_1} 2 \qquad\qquad \text{or} \qquad\qquad
\begin{tikzcd}[row sep=small, column sep=small]
 & 0 \arrow[dl] \\
1 & \\
 & 2 \arrow[ul]
\end{tikzcd}
\qquad .
\]
They correspond respectively to the diagrams
\[
a_0 \xleftarrow{u_0} a_1 \xrightarrow{u_1} a_2
\]
and
\[
a_0 \longrightarrow a_1 \longleftarrow a_2.
\]
The categories $\Phi_1$, $\Psi_1$ are opposite to one another, and are not isomorphic. So there is in general no isomorphism between $\Ch_{\Phi_1}(X)$ and $\Ch_{\Psi_1}(X)$.
\end{enumerate}

Let $I$ be an oriented interval, and let us consider the \textit{opposite interval $I^{\opp}$}, having the same underlying category but the opposite orientation. As the definition of the category of paths of type $I$ does not involve the orientation of $I$, we thus have a canonical isomorphism % [sic] l'orientation opposé -> l'orientation opposée, agreement
\[
\Ch_I(X) \xrightarrow{\sim} \Ch_{I^{\opp}}(X).
\]

\msp{page 17}

Nevertheless we shall take care not to confuse a path of type $I$ with the opposite path of type $I^{\opp}$, the source of the one being the target of the other and conversely. This is to say that a path must be seen as \textit{more} than a mere functor
\[
I \longrightarrow X,
\]
we must in addition give ourselves the \textit{orientation of $I$} as forming part of the path structure. So the reasonable formal definition, which fits the intuition of a path as going from a well-defined \textit{source} toward an equally well-defined \textit{target}, is that a path (defined in terms of an interval category $I$) is the datum of a pair $(\omega,c)$, where $\omega$ is an orientation of $I$, and $c$ a functor $c : I \longrightarrow X$. When we look at the category of paths of a fixed ``type'', this \textit{type} is $(I,\omega)$, and not $I$. It is when $\omega$ is fixed once and for all that this category can with impunity be identified with $\uHom(I,X)$ without more. % [sic] la donné -> la donnée, missing final e

We must view this category of paths as equipped with the essential structure
\[
\begin{tikzcd}
\Ch_{(I,\omega)}(X) \arrow[d, "{(\sigma,\beta)}"] \\
X \times X
\end{tikzcd}
\qquad
\begin{array}{l}
c \longmapsto (\sigma(c),\beta(c)) \\
\text{where $\sigma(c)$ is the \textit{source,}} \\
\qquad \text{$\beta(c)$ the \textit{target} of $c$.}
\end{array} % [sic] ou sigma(c) -> où sigma(c), missing accent
\]

\printedpage{12}
% head: A. Grothendieck : Dérivateurs, Ch. VII, Édition des Universités de Montpellier II et Paris VII
% foot: Version 22 décembre 2001

\msp{page 18}

This functor is identified with the composition functor
\[
c \longmapsto c \circ \alpha
\]
where $\alpha$ is the inclusion functor
\[
\left\{
\begin{array}{lll}
\varepsilon \xrightarrow{\alpha} I & & \\[1ex]
\varepsilon = \{0,1\}, & \multicolumn{2}{l}{\text{discrete category with two elements,}} \\[1ex]
\alpha \text{ given by} & \alpha(0) = & \text{source of } I \\[1ex]
 & \alpha(1) = & \text{target of } I.
\end{array}
\right.
\]
If $\omega'$ is the opposite orientation of $I$, we have a commutative diagram
\[
\begin{tikzcd}
\Ch_{(I,\omega)}(X) \arrow[r, "\varphi", "\sim"'] \arrow[d, "{(\sigma,\beta)}"] & \Ch_{(I,\omega')}(X) \arrow[d, "{(\sigma',\beta')}"] \\
X \times X \arrow[r, "s"] & X \times X,
\end{tikzcd}
\]
where $s$ is the symmetry functor
\[
s(x,y) = (y,x)
\]
(\footnote[6]{The functor $\varphi$ is the one which associates to every path $c$, of type $\tau$, the \textit{opposite path $c^{\opp}$} or $c^{-1}$, of type $\tau^{\opp}$. If $\tau = \tau^{\opp}$, this is therefore a path \textit{of the same type,} and $\varphi$ is an \textit{automorphism} of $\Ch_\tau(X)$. This automorphism (functorial in $X$) is deduced from an automorphism $\theta$ of $I$, cf. below.}). Thus the isomorphism $\varphi$ does not respect the structural ``source-target'' arrows $(\sigma,\beta)$.

There is however a difficulty when there exists an isomorphism of \textit{oriented} intervals
\[
\theta : (I,\omega) \xrightarrow{\sim} (I,\omega'),
\]
i.e. when $I$ is symmetric. This is not the case for the $\bDelta_i$ $(i \geq 1)$ (these categories have no automorphisms other than the identity), but it is the case for the categories $\Phi_1$, $\Psi_1$ of example $4^o$, more generally \msp{page 19} for the reduced categories whose length is \textit{even.} As every time we have an isomorphism $\theta$ of intervals that are \textit{oriented,} we find an isomorphism between the categories of paths of type $(I,\omega)$ and $(I,\omega)^{\opp} = (I,\omega')$, \textit{compatible} this time with the functors $(\sigma,\beta)$:
\[
\begin{tikzcd}
\Ch_{(I,\omega)}(X) \arrow[rr, "\Psi_\theta", "\sim"'] \arrow[dr, "{(\sigma,\beta)}"'] & & \Ch_{(I,\omega')}(X) \arrow[dl, "{(\sigma',\beta')}"] \\
 & X \times X. &
\end{tikzcd}
\]

\printedpage{13}
% head: A. Grothendieck : Dérivateurs, Ch. VII, Édition des Universités de Montpellier II et Paris VII
% foot: Version 22 décembre 2001

\msp{page 20}

It is good to regard a \textit{type} of paths as a noncategorical, but combinatorial, datum, a simple sequence $\tau$ of $n$ signs $\tau_i$ $(1 \leq i \leq n)$ $+$ or $-$ ($n =$ length of the paths under consideration). This combinatorial datum defines a standard oriented interval, $I_\tau$. When we have an arbitrary oriented interval $(I,\omega)$, it defines a type $\tau$, and we identify it canonically with the type $I_\tau$. It is better to write $\Ch_\tau(X)$, rather than $\Ch_{(I,\omega)}(X)$, to make it quite clear that for given $X$, this is a totally ``pinned down'' category, without automorphisms tied to the particular nature of $\tau$ (to automorphisms of $\tau$!). % [sic] une donné -> une donnée, missing final e; [sic] Cette donné -> Cette donnée, missing final e

We can make this point precise as follows. Let $\tau'$ be another type of intervals, and let us set out to determine ``laws'' which allow us to associate to any path of type $\tau$ in an arbitrary category $X$, a path of type $\tau'$ of this same category, and this in a ``functorial'' way, i.e. compatible with the horizontal maps
\[
\begin{tikzcd}
\mathrm{Ch}_\tau(X) \arrow[r, "\mathrm{Ch}_\tau(f)"] \arrow[d, "L_X"'] & \mathrm{Ch}_\tau(X') \arrow[d, "L_{X'}"] \\
\mathrm{Ch}_{\tau'}(X) \arrow[r, "\mathrm{Ch}_{\tau'}(f)"] & \mathrm{Ch}_{\tau'}(X')
\end{tikzcd}
\]
associated with a functor $f : X \longrightarrow X'$.

\msp{page 21}

For example, taking the type $\bDelta_2^+$ (i.e. $+1,+1$), corresponding to the paths
\[
a_0 \xrightarrow{u_1} a_1 \xrightarrow{u_2} a_2,
\]
we can associate to such a path the path of type \texttt{[}$\bDelta_1^+$\texttt{]}
\[
a_0 \xrightarrow{u_2u_1} a_2,
\]
or also the path of type $\bDelta_2^+$
\[
a_0 \xrightarrow{u_2u_1} a_2 \xrightarrow{\mathrm{id}} a_2,
\]
etc. Since the functor
\[
\left\{
\begin{array}{cccl}
X & \longmapsto & \mathrm{Ch}_\tau(X) & \qquad (= \mathrm{Ob}\, \Ch_\tau(X)) \\[1ex]
\mathrm{Cat} & \longrightarrow & (\mathrm{Ens}) &
\end{array}
\right.
\]
is \textit{represented} by the object $I_\tau$ of $X$ \texttt{[rather of} Cat\texttt{]}, the sorites of representable functors tells us that such laws of formation of paths of type $\tau$ from paths of type $\tau'$ \texttt{[rather} $\tau'$ \texttt{from paths of type} $\tau$\texttt{]} correspond bijectively to the arrows in Cat
\[
I_{\tau'} \xrightarrow{\varphi} I_\tau\ ,
\]

\end{document}
